
\documentclass[conference]{IEEEtran}

\usepackage{graphicx}

\ifCLASSINFOpdf
  
\else
 
\fi


\hyphenation{}


\begin{document}

\title{[Placeholder Method] for [PlaceHolder Goal Achieved] for [PlaceHolder Application]}

\author{\IEEEauthorblockN{[PlaceHolder Name]}
\IEEEauthorblockA{Deparment of [PlaceHolder] Engineering \\
University of [PlaceHolder]\\
PlaceHolder Address\\
Email: [PlaceHolder]}
\and
\IEEEauthorblockN{[PlaceHolder Name]}
\IEEEauthorblockA{Deparment of [PlaceHolder] Engineering \\
University of [PlaceHolder]\\
PlaceHolder Address\\
Email: [PlaceHolder]}
}

% conference papers do not typically use \thanks and this command
% is locked out in conference mode. If really needed, such as for
% the acknowledgment of grants, issue a \IEEEoverridecommandlockouts
% after \documentclass

% for over three affiliations, or if they all won't fit within the width
% of the page, use this alternative format:
% 
%\author{\IEEEauthorblockN{Michael Shell\IEEEauthorrefmark{1},
%Homer Simpson\IEEEauthorrefmark{2},
%James Kirk\IEEEauthorrefmark{3}, 
%Montgomery Scott\IEEEauthorrefmark{3} and
%Eldon Tyrell\IEEEauthorrefmark{4}}
%\IEEEauthorblockA{\IEEEauthorrefmark{1}School of Electrical and Computer Engineering\\
%Georgia Institute of Technology,
%Atlanta, Georgia 30332--0250\\ Email: see http://www.michaelshell.org/contact.html}
%\IEEEauthorblockA{\IEEEauthorrefmark{2}Twentieth Century Fox, Springfield, USA\\
%Email: homer@thesimpsons.com}
%\IEEEauthorblockA{\IEEEauthorrefmark{3}Starfleet Academy, San Francisco, California 96678-2391\\
%Telephone: (800) 555--1212, Fax: (888) 555--1212}
%\IEEEauthorblockA{\IEEEauthorrefmark{4}Tyrell Inc., 123 Replicant Street, Los Angeles, California 90210--4321}}




% use for special paper notices
%\IEEEspecialpapernotice{(Invited Paper)}




% make the title area
\maketitle

% As a general rule, do not put math, special symbols or citations
% in the abstract
\begin{abstract}
[PlaceHolder] is a huge problem. Statistics of how huge problem [PlaceHolder]. Traditionally people solve it with [PlaceHolder]. However, [PlaceHolder] has not shown much improvement for [PlaceHolder] application, which this paper proposes improvement in [PlaceHolder] metrics with utilization of [PlaceHolder] technology compared to [PlaceHolder] traditional technology.
\end{abstract}

% no keywords




% For peer review papers, you can put extra information on the cover
% page as needed:
% \ifCLASSOPTIONpeerreview
% \begin{center} \bfseries EDICS Category: 3-BBND \end{center}
% \fi
%
% For peerreview papers, this IEEEtran command inserts a page break and
% creates the second title. It will be ignored for other modes.
\IEEEpeerreviewmaketitle




\section{Introduction}
~\textbf{This is a huge problem with absurd challenges} Human interaction increases demand increases, rise of wireless device communication challenges increase.
\\
~\textbf{Traditional methods can not handle this much dynamics} Hard handshake was there, then soft handshake got introduced. Still they could not handle channel dynamics. Important is to learn the adversaries of the channel being utilized and when it is available to facilitate the secondary user.
\\
~\textbf{How people utilized unsupervised, SNN to resolve these problems}
\\
~\textbf{Here's our contribution:}
1. Dynamic Channel Condition
2. Unsupervised Learning
3. Robust model, real life simulation with Raspberry pi
\section{Methodology}
Overview of the figure~\ref{fig:methodology}. The input data is Dynamic channel path condition. Output is detecting availibility in primary users.
\begin{figure}
    \centering
    \includegraphics[width=\linewidth]{traffic_flowchart.png}
    \caption{Overview of the methodology}
    \label{fig:methodology}
\end{figure}

Centralized learning works because Secondary Users (SUs) are doing individual training of their own data.



% conference papers do not normally have an appendix


% use section* for acknowledgment
\section*{Acknowledgment}


The authors would like to thank...





% trigger a \newpage just before the given reference
% number - used to balance the columns on the last page
% adjust value as needed - may need to be readjusted if
% the document is modified later
%\IEEEtriggeratref{8}
% The "triggered" command can be changed if desired:
%\IEEEtriggercmd{\enlargethispage{-5in}}

% references section

% can use a bibliography generated by BibTeX as a .bbl file
% BibTeX documentation can be easily obtained at:
% http://mirror.ctan.org/biblio/bibtex/contrib/doc/
% The IEEEtran BibTeX style support page is at:
% http://www.michaelshell.org/tex/ieeetran/bibtex/
%\bibliographystyle{IEEEtran}
% argument is your BibTeX string definitions and bibliography database(s)
%\bibliography{IEEEabrv,../bib/paper}
%
% <OR> manually copy in the resultant .bbl file
% set second argument of \begin to the number of references
% (used to reserve space for the reference number labels box)
\begin{thebibliography}{1}

\bibitem{IEEEhowto:kopka}
H.~Kopka and P.~W. Daly, \emph{A Guide to \LaTeX}, 3rd~ed.\hskip 1em plus
  0.5em minus 0.4em\relax Harlow, England: Addison-Wesley, 1999.

\end{thebibliography}




% that's all folks
\end{document}


